\mxnewpage\section*{EV4-B -- Préparation EAM}\label{EV4-B}
\addcontentsline{toc}{section}{EV4-B -- Préparation EAM}
\textbf{Sujet réservé au professeur}. Ensemble du programme.\par Durée : 2 h. Sans calculatrice. Note sur 20. Entraînements de mai et juin ; EAM de zone à confirmer.
\subsection*{Partie 1 -- Automatismes (6 points)}
Douze questions, \num{0,5} point chacune. Une seule réponse correcte ; aucune justification demandée. Une réponse incorrecte ou absente rapporte 0, sans retrait de points.
\begin{enumerate}
\item Calculer $3+2\times3$.\par A.~$9$\quad B.~$12$\quad C.~$15$\quad D.~$8$
\item Calculer $\frac12+\frac13$.\par A.~$\frac{5}{3}$\quad B.~$\frac{5}{6}$\quad C.~$\frac{2}{5}$\quad D.~$\frac{1}{6}$
\item Calculer $\dfrac{2^3}{2^{1}}$.\par A.~$8$\quad B.~$2$\quad C.~$4$\quad D.~$1$
\item Une valeur de 50 augmente de \pc{20}. Quelle est la valeur finale ?\par A.~$70$\quad B.~$40$\quad C.~$60$\quad D.~$10$
\item Une valeur de 80 est obtenue après une baisse de \pc{20}. Quelle est la valeur initiale ?\par A.~$60$\quad B.~$64$\quad C.~$96$\quad D.~$100$
\item Quel est le coefficient multiplicateur global d’une hausse de \pc{25} suivie d’une baisse de \pc{20} ?\par A.~$\frac{21}{20}$\quad B.~$1$\quad C.~$\frac{4}{5}$\quad D.~$\frac{9}{10}$
\item La solution de l’équation $2x+3=0$ est~:\par A.~$\frac{3}{2}$\quad B.~$6$\quad C.~$\frac{-3}{2}$\quad D.~$-6$
\item La pente de la droite passant par $\pt{0;3}$ et $\pt{2;9}$ est~:\par A.~$6$\quad B.~$-3$\quad C.~$3$\quad D.~$2$
\item La moyenne de $3, 5, 7$ est~:\par A.~$3$\quad B.~$15$\quad C.~$6$\quad D.~$5$
\item Un secteur de $90^\circ$ représente quelle fraction du total ?\par A.~$\frac{1}{3}$\quad B.~$\frac{1}{2}$\quad C.~$\frac{1}{4}$\quad D.~$\frac{3}{4}$
\item Une probabilité vaut $\dfrac{2}{5}$. Son complément est~:\par A.~$\frac{3}{5}$\quad B.~$\frac{1}{5}$\quad C.~$\frac{2}{5}$\quad D.~$\frac{5}{2}$
\item Quelle est la valeur de $f(-2)$ pour $f(x)=x^2+3$ ?\par A.~$12$\quad B.~$5$\quad C.~$-1$\quad D.~$7$
\end{enumerate}
\mxnewpage\subsection*{Partie 2 -- Exercices indépendants (14 points)}
Toute démarche doit être justifiée. Les résultats exacts sont attendus sauf indication contraire.
\mxneed{12}
\subsection*{Exercice 1 -- Un tarif moyen (5 points)}
Un service propose plusieurs tarifs. 60 pour cent des clients choisissent la formule A ; parmi eux, la moitié paie 0 unité et l’autre moitié paie 4 unités. Tous les autres clients paient 6 unités. On choisit un client au hasard et l’on note $X$ le prix qu’il paie. La dernière question étudie une réduction uniforme, y compris pour le tarif initialement nul.
\begin{enumerate}
\item Établir la loi de probabilité de $X$. \hfill(1 pt)
\item Calculer son espérance et sa variance. \hfill(2 pts)
\item Tous les prix diminuent de 1 unité. Calculer la nouvelle moyenne et le nouvel écart type. Interpréter le prix négatif obtenu pour certains clients. \hfill(2 pts)
\end{enumerate}
\mxneed{12}
\subsection*{Exercice 2 -- Maximum exponentiel (5 points)}
La fonction $f$ est définie sur $\intff{0}{4}$ par $f(t)=t^2e^{-t}$. On cherche son maximum sur cet intervalle. Pour la dernière comparaison sans calculatrice, on pourra utiliser l’inégalité donnée $e^2>4$.
\begin{enumerate}
\item Calculer $f\prime$ sous forme factorisée. \hfill(1 pt)
\item Déterminer les variations et le maximum. \hfill(2 pts)
\item Comparer exactement $f(4)$ au maximum, sachant $e^2>4$. \hfill(2 pts)
\end{enumerate}
\mxneed{12}
\subsection*{Exercice 3 -- Déterminer une position (4 points)}
Dans un repère orthonormé, les points $A\pt{0;0}$ et $B\pt{4;0}$ sont fixes. Le point $M$ se déplace sur la droite d’équation $y=2$ ; ses coordonnées sont donc $M\pt{x;2}$, où $x$ parcourt $\R$.
\begin{enumerate}
\item Calculer $\vect{MA}\cdot\vect{MB}$. \hfill(1 pt)
\item Déterminer les valeurs de $x$ pour lesquelles les vecteurs $\vect{MA}$ et $\vect{MB}$ sont orthogonaux. \hfill(1 pt)
\item Interpréter le résultat à l’aide du cercle de diamètre AB. \hfill(2 pts)
\end{enumerate}
